\documentclass[10pt,letterpaper]{article}
\usepackage{cornell-notes}

\title{Clause 24.03.07.07: Tuple Interface}
\author{}
\date{}
\setDocTitle{Clause 24.03.07.07: Tuple Interface}
\setDocSubtitle{C++ 2024}
\setDocOwner{C++ 2024 Cornell Notes}
\setDocDate{\today}
\setCornellCollection{C++ 2024 Cornell Notes}
\setCornellUnitType{Clause}
\setCornellUnitNumber{24.03.07.07}
\setCornellUnitTitle{Tuple Interface}
\hypersetup{pdftitle={Clause 24.03.07.07: Tuple Interface},pdfsubject={Cornell notes},pdfauthor={}}

\begin{document}
\maketitle
\begin{CornellOverviewBox}{Overview}
\textbf{Classification:} Standard library - Normative\\[2pt]
\textbf{Learning objective:} Understand the rules, terminology, and semantic consequences associated with Tuple interface.
\end{CornellOverviewBox}\begin{CornellSummaryBox}{Summary}
{\sffamily\bfseries\color{Accent} Summary}\par\vspace{3pt}Section 24.3.7.7 focuses on Tuple interface. Apply the specified interface together with its mandates, effects, result conditions, exception behavior, and complexity constraints. When applying it, preserve the distinction between mandatory requirements, recommendations, permissions, and behavior deliberately left undefined, unspecified, or implementation-defined.
\end{CornellSummaryBox}

\end{document}
